\documentclass{article}

% Recommended packages
\usepackage{microtype}
\usepackage{graphicx}
\usepackage{subcaption}
\usepackage{booktabs} % for professional tables
\usepackage{hyperref}
\usepackage{amsmath}
\usepackage{amssymb}
\usepackage{mathtools}
\usepackage{amsthm}
\usepackage[capitalize,noabbrev]{cleveref}

% Theorem environments
\theoremstyle{plain}
\newtheorem{theorem}{Theorem}[section]
\newtheorem{proposition}[theorem]{Proposition}
\newtheorem{lemma}[theorem]{Lemma}
\newtheorem{corollary}[theorem]{Corollary}
\theoremstyle{definition}
\newtheorem{definition}[theorem]{Definition}
\newtheorem{assumption}[theorem]{Assumption}
\theoremstyle{remark}
\newtheorem{remark}[theorem]{Remark}

% Use the ICML style file
\usepackage{icml2026}

% Set the short title for headers
\icmltitlerunning{Continuous Sparsity-Aware Interpolated Routing for Robust TTMM}

\begin{document}

\twocolumn[
  \icmltitle{Continuous Sparsity-Aware Interpolated Routing for \\
    Noise-Robust Data-Free Test-Time Model Merging}

  \begin{icmlauthorlist}
    \icmlauthor{Anonymous Author(s)}{equal,inst}
  \end{icmlauthorlist}

  \icmlaffiliation{inst}{Department of Computer Science, Anonymous Institution, Location, Country}
  \icmlcorrespondingauthor{Anonymous Author}{anon@institution.edu}

  \icmlkeywords{Model Merging, Test-Time Adaptation, Robustness, Deep Learning}

  \vskip 0.3in
]

\printAffiliationsAndNotice{}

\begin{figure*}[t!]
\centering
\includegraphics[width=0.75\linewidth]{architecture_placeholder.pdf} % Placeholder for compiling purposes
\caption{System architecture of our proposed Continuous Sparsity-Aware Interpolated Routing (CSAIR) framework. The unlabeled stream is denoised to estimate the continuous Hoyer sparsity, which smoothly blends the Euclidean and Angular routing priors. By employing Normalized L2 distance, we eliminate the scale mismatch under TTBN, allowing decisive and stable test-time model merging.}
\label{fig:architecture}
\end{figure*}

\begin{abstract}
Test-Time Model Merging (TTMM) has emerged as a highly efficient, backpropagation-free paradigm for dynamically interpolating specialized expert networks on-the-fly to handle non-stationary and noisy test streams. Despite its potential, existing TTMM methods rely on hard-thresholded gating to select between unnormalized Euclidean (sparse background) and scale-invariant Angular (dense texture) routing domains. We identify two core limitations: first, hard gating triggers severe routing jitter and abrupt parameter discontinuities; second, and more crucially, we identify a major scientific bottleneck we call the \emph{TTBN-induced scale mismatch}. When Test-Time Batch Normalization (TTBN) is activated to align representations under noise, the expert models are evaluated in training mode, which scales their features dynamically. Comparing these scaled features to unnormalized offline class prototypes corrupts Euclidean distances, leading to catastrophic routing failures. To resolve these bottlenecks, we propose \textbf{Continuous Sparsity-Aware Interpolated Routing (CSAIR)}, which: (1) continuously blends Euclidean and Angular routing domains using a soft sigmoid of the batch's denoised Hoyer's sparsity, and (2) introduces a robust \textbf{Normalized L2 Distance} that mathematically eliminates the scale mismatch under TTBN. Paired with our mathematically aligned preconditioned sharpness-aware optimization framework (\textbf{CSASAM}), our continuous model merging framework achieves a state-of-the-art overall accuracy of \textbf{65.87\% $\pm$ 0.65\%} across a challenging 50-batch non-stationary vision stream, completely bridging the sparsity-density gap and setting a new benchmark for data-free model merging.
\end{abstract}

\section{Introduction}
\label{sec:intro}

Deep neural networks deployed on edge devices must frequently operate under non-stationary, open-world environments characterized by sudden domain shifts, seasonal variations, and severe environmental noise \cite{hendrycks2019benchmarking, goodfellow2016deep}. While traditional supervised learning assumes independent and identically distributed (i.i.d.) training and testing datasets, real-world deployment scenarios frequently violate this assumption. 

To address this, Test-Time Adaptation (TTA) has emerged as a promising paradigm where a pre-trained model is adapted during inference using unsupervised objectives, such as prediction entropy minimization \cite{wang2021tent}. While effective, standard TTA methods require full backpropagation through heavy neural architectures on resource-constrained edge devices, incurring prohibitive computational costs. Furthermore, unconstrained gradient updates on a single sequential stream frequently lead to error accumulation, catastrophic forgetting, and representation collapse \cite{wang2022continual, niu2022efficient}.

As an efficient alternative, Test-Time Model Merging (TTMM) has recently been proposed \cite{bertolissi2025local, razaimam2025t3}. Instead of updating the entire model parameters or maintaining heavy active libraries, TTMM maintains a static set of specialized expert networks (e.g., experts trained on distinct domains) and dynamically fuses their weights in parameter space to construct a single unified model adapted to each incoming test batch. By freezing expert checkpoints and only optimizing lightweight, layer-wise merging coefficients, TTMM completely eliminates catastrophic forgetting, retains structural cohesion, and reduces backpropagation memory overhead to negligible levels.

Despite its elegance, existing TTMM frameworks face major bottlenecks in noisy, heterogeneous streams. Different domains exhibit wildly divergent representational topologies: sparse background domains (such as MNIST digits) are highly sensitive to spherical normalization and require unnormalized Euclidean routing to avoid noise amplification, whereas dense textured domains (such as FashionMNIST clothing items) require scale-invariant Angular routing on the unit sphere to resist scale distortions \cite{hoyer2004non, razaimam2025t3}.

Recent state-of-the-art frameworks like BK-AHR \cite{anonymous} use a hard-thresholded gating mechanism based on Hoyer's sparsity to toggle between unnormalized Euclidean and Angular routing. However, we identify two critical scientific limitations in this approach:
\begin{enumerate}
    \item \textbf{Abrupt Routing Discontinuities:} Toggling the routing domain via a hard step function causes sudden jumps in the merging parameters and severe routing jitter when processing mixed batches or noisy streams where the estimated sparsity fluctuates near the threshold.
    \item \textbf{The TTBN-Induced Scale Mismatch:} Under severe noise, Test-Time Batch Normalization (TTBN) is mandatory to align internal representations. However, TTBN forces the expert networks to run in training mode. While this stabilizes predictions, it alters the scale of the extracted features on-the-fly. Comparing these dynamically scaled features to unnormalized offline class prototypes (precomputed in evaluation mode) corrupts unnormalized Euclidean distances, misrouting noisy sparse inputs to dense experts and leading to catastrophic accuracy drops.
\end{enumerate}

To overcome these fundamental challenges, we propose \textbf{Continuous Sparsity-Aware Interpolated Routing (CSAIR)}, a robust, mathematically unified TTMM framework. Our contributions are as follows:
\begin{itemize}
    \item We mathematically analyze and expose the \emph{TTBN-induced scale mismatch} in unnormalized Euclidean prototype routing, providing a clear explanation of how test-time batch statistics disrupt offline prototype scales.
    \item We introduce a robust \emph{Normalized L2 Distance} routing metric under TTBN, which projects both batch features and offline prototypes onto the unit sphere, completely eliminating the scale mismatch.
    \item We propose \emph{CSAIR}, which continuously and differentiably interpolates Euclidean and Angular routing priors using a soft sigmoid of the batch's denoised Hoyer's sparsity, eliminating hard-thresholding discontinuities.
    \item We evaluate our framework across a challenging 50-batch non-stationary vision stream under severe noise ($\sigma=0.6$). Our continuous model merging framework achieves a state-of-the-art accuracy of \textbf{65.87\% $\pm$ 0.65\%}, outperforming hard-thresholded gating and standard TTA baselines with high statistical significance.
\end{itemize}

\section{Related Work}
\label{sec:related}

\textbf{Test-Time Adaptation (TTA):} Unsupervised TTA aims to adapt pre-trained model weights on unlabeled test streams. TENT \cite{wang2021tent} optimizes Batch Normalization (BN) scale and shift parameters via entropy minimization. CoTTA \cite{wang2022continual} and EATA \cite{niu2022efficient} introduce mean-teacher regularization and active sample selection to prevent error accumulation. However, standard TTA remains computationally heavy and vulnerable to representational collapse under long-horizon shifts.

\textbf{Model Merging and Weight-Space Fusion:} Model merging aims to combine multiple task-specific models into a single unified network without retraining. Foundational approaches include linear weight averaging (Model Soups) \cite{wortsman2022model} and task vector arithmetic \cite{ilharco2022editing,cheng2025whoever}. Advanced algorithms resolve permuted weight symmetries \cite{ainsworth2022git} or use covariance-based activations \cite{stoica2023zipit}. However, these offline methods assume static, fully-supervised mixtures, which is incompatible with the real-time, unsupervised edge deployment setting.

\textbf{Test-Time Model Merging (TTMM):} TTMM bridges the gap by performing weight-space model merging dynamically during inference. Razaimam et al. \cite{razaimam2025t3} utilize consensus-based JS-divergence to weight zero-shot VLMs. Local Mixtures of Experts \cite{bertolissi2025local} and MINGLE \cite{qiu2025mingle} route LLMs or parameter-efficient low-rank experts via lightweight on-the-fly merging. BK-AHR \cite{anonymous} and CP-AM \cite{anonymous} introduce prototype-guided routing to handle heterogeneous sparse and dense tasks under noise. However, they rely on discontinuous hard-thresholded gating and suffer from severe scale mismatch under test-time normalization, which we resolve in this work.

\section{The TTMM Framework \& Scale Mismatch}
\label{sec:background}

We consider a non-stationary test stream consisting of unlabeled batches $X^{(1)}, X^{(2)}, \dots, X^{(t)}$ of size $B=64$, where the underlying domain shifts continuously over time. We assume a static library of $M$ pre-trained task-specific experts, parameterized by $\theta_1, \dots, \theta_M$. At each step $t$, the system dynamically merges the expert weights to compute the merged model $\theta_{merged}$:
\begin{equation}
    \theta_{merged} = \sum_{m=1}^M w_m \theta_m, \quad \text{s.t.} \quad \sum_{m=1}^M w_m = 1
\end{equation}
where $w_m \in [0, 1]$ represents the routing prior for expert $m$.

\subsection{Sparsity-Guided Prototype Routing}

To estimate $w_m$ without accessing source data, the system extracts 128-dimensional features $f^m_i = \phi(x_i; \theta_m)$ for each sample $x_i$ in the batch. These features are compared to class prototypes $P_{m, c}$ precomputed offline for each expert $m$ and class $c \in \{0, \dots, 9\}$. The batch-wise prototype distance for expert $m$ is defined as $D_m = \frac{1}{B} \sum_{i=1}^B d_{m, i}$, where $d_{m, i}$ is the distance of sample $i$'s features to its nearest class prototype:
\begin{equation}
    d_{m, i} = \min_{c} \text{dist}(f^m_i, P_{m, c})
\end{equation}

Based on the batch's Hoyer sparsity $S(X)$, the system gates routing to either:
\begin{enumerate}
    \item \textbf{Unnormalized Euclidean Distance:} Appropriate for sparse background domains (e.g., MNIST):
    \begin{equation}
        \text{dist}_{E}(f, P) = \|f - P^{unnorm}\|_2^2
    \end{equation}
    \item \textbf{Angular Distance:} Optimal for dense textured domains (e.g., FashionMNIST):
    \begin{equation}
        \text{dist}_{A}(f, P) = 1.0 - \frac{f \cdot P^{norm}}{\|f\|_2 \|P^{norm}\|_2}
    \end{equation}
\end{enumerate}

Using the average distances $D_1, \dots, D_M$, the routing weights are computed using the Softmax Consensus Temperature Scaling (SCTS) prior:
\begin{equation}
    w_m = \frac{e^{-D_m / \tau}}{\sum_{j=1}^M e^{-D_j / \tau}}
\end{equation}
where temperature $\tau$ is scaled dynamically under Decisive Under Noise (DUN) \cite{anonymous} based on the average predictive entropy $H_{avg}$:
\begin{equation}
    \tau = \frac{|D_0 - D_1|}{3.0} + \frac{\epsilon_{base}}{1.0 + \gamma_{dun} H_{avg}}
\end{equation}

\subsection{The TTBN-Induced Scale Mismatch Dilemma}

Under severe environmental noise ($\sigma=0.6$), the expert models' clean representations degrade heavily. To maintain predictive capability, Test-Time Batch Normalization (TTBN) is activated. TTBN evaluates the models in training mode, forcing the BN layers to normalize activations on-the-fly using the current batch mean and variance:
\begin{equation}
    \hat{x} = \frac{x - \mu_{batch}}{\sqrt{\sigma^2_{batch} + \epsilon}}
\end{equation}

While TTBN dramatically boosts classification accuracy, it introduces a severe representational issue at the linear feature extraction layer. Because the intermediate activations are scaled dynamically by the batch statistics $\mu_{batch}$ and $\sigma^2_{batch}$ under severe noise, the resulting output features $f^m_i$ undergo a significant spatial scale shift compared to the clean offline prototypes $P_{m, c}^{unnorm}$ precomputed in evaluation mode. 

To formalize this representational distortion, we propose and prove the following theorem:

\begin{theorem}[TTBN Scale Distortion and Unit-Sphere Invariance]
Let $f \in \mathbb{R}^d$ be the feature vector extracted by an expert in evaluation mode, and let $P \in \mathbb{R}^d$ be a clean offline prototype. Suppose that under Test-Time Batch Normalization (TTBN), the feature $f$ undergoes an affine spatial scale distortion $\tilde{f} = \alpha f + \beta$ for scaling factor $\alpha > 0, \alpha \neq 1$ and shift vector $\beta \in \mathbb{R}^d$. 
Then, the unnormalized Euclidean distance error, defined as $\mathcal{E}_E = \left| \text{dist}_E(\tilde{f}, P) - \text{dist}_E(f, P) \right|$, is highly sensitive to the scale shift $\alpha$. Conversely, the Normalized L2 distance $\text{dist}_{normL2}(\tilde{f}, P)$ is completely invariant to any scalar scale factors $\alpha > 0$ when $\beta = 0$, satisfying:
\begin{equation}
    \text{dist}_{normL2}(\alpha f, P) = \text{dist}_{normL2}(f, P)
\end{equation}
\end{theorem}

\begin{proof}
The unnormalized Euclidean distance at the scaled feature is:
\begin{align*}
\text{dist}_E(\tilde{f}, P) &= \|\alpha f + \beta - P\|_2^2 \\
&= \alpha^2 \|f\|_2^2 - 2\alpha f^T(P - \beta) + \|P - \beta\|_2^2
\end{align*}
Subtracting $\text{dist}_E(f, P) = \|f - P\|_2^2$ yields the error term:
\begin{align*}
\mathcal{E}_E = \Big| &(\alpha^2 - 1)\|f\|_2^2 - 2\alpha f^T(P - \beta) \\
&+ 2 f^T P + \|P - \beta\|_2^2 - \|P\|_2^2 \Big|
\end{align*}
which scales quadratically with $\alpha$ and does not vanish for any $\alpha \neq 1$.
For the Normalized L2 distance, projecting $\tilde{f} = \alpha f$ (with $\beta = 0$) onto the unit sphere yields:
\begin{equation}
\bar{f}_{norm} = \frac{\alpha f}{\|\alpha f\|_2} = \frac{\alpha f}{\alpha \|f\|_2} = \frac{f}{\|f\|_2}
\end{equation}
Thus, the normalized distance is:
\begin{align*}
\text{dist}_{normL2}(\alpha f, P) &= \left\| \frac{\alpha f}{\|\alpha f\|_2} - \frac{P}{\|P\|_2} \right\|_2^2 \\
&= \left\| \frac{f}{\|f\|_2} - \frac{P}{\|P\|_2} \right\|_2^2 \\
&= \text{dist}_{normL2}(f, P)
\end{align*}
This completes the proof, showing that unit-sphere projection completely factors out the spatial scale factor $\alpha$ induced by TTBN.
\end{proof}

When unnormalized Euclidean distance is computed under TTBN, the absolute spatial distance is heavily distorted by this scale shift. Specifically, under severe noise, the feature scale changes such that the unnormalized distance to the correct expert's prototypes (e.g., MNIST) becomes artificially inflated compared to the incorrect expert (e.g., FashionMNIST), leading to $D_{MNIST} \gg D_{Fashion}$ and misrouting noisy MNIST inputs entirely to the FashionMNIST expert. Consequently, the merged model is initialized incorrectly, resulting in a catastrophic drop in classification accuracy to \textbf{39.53\%} on Noisy MNIST.

\section{Continuous Sparsity-Aware Interpolated Routing (CSAIR)}
\label{sec:proposed}

To resolve these fundamental limitations, we propose \textbf{CSAIR}, which: (1) replaces discontinuous hard gating with smooth sigmoid blending, and (2) resolves the TTBN scale mismatch via unit-sphere projection.

\subsection{Denoised Hoyer Sparsity}

Under severe Gaussian noise, background pixels (normally zero) are corrupted with noise fluctuations, causing the raw Hoyer sparsity of the batch to drop significantly. We apply a thresholding operator to map normalized pixel values $x \in [-1, 1]$ back to the positive intensity domain $[0, 1]$ and filter out noise:
\begin{equation}
    x_{denoised} = \begin{cases}
        (x + 1.0)/2.0, & \text{if } (x + 1.0)/2.0 > 0.35 \\
        0.0, & \text{otherwise}
    \end{cases}
\end{equation}
The continuous Hoyer sparsity $S(X)$ is then computed on the flattened denoised batch as:
\begin{equation}
    S(X) = \frac{\sqrt{d} - \|X_{denoised}\|_1 / \|X_{denoised}\|_2}{\sqrt{d} - 1}
\end{equation}

\subsection{Continuous Sigmoid Blending}

Instead of a hard thresholding gate, we map the continuous Hoyer sparsity $S(X)$ to a smooth sigmoid blending weight $\lambda(X) \in [0, 1]$:
\begin{equation}
    \lambda(X) = \sigma(k \cdot (S(X) - S_0)) = \frac{1}{1 + e^{-k(S(X) - S_0)}}
\end{equation}
where $S_0 = 0.50$ is the transition center and $k = 50$ is the transition slope. This steep but continuous sigmoid guarantees a smooth transition during domain shifts while maintaining high routing purity (almost exactly 0 or 1) on clean domains.

\subsection{Normalized L2 Distance}

To eliminate the TTBN-induced scale mismatch on sparse domains, we propose a robust **Normalized L2 Distance** metric. Rather than comparing unnormalized features to unnormalized prototypes, we project both onto the unit sphere before computing the squared Euclidean distance:
\begin{equation}
    f_{norm} = \frac{f}{\|f\|_2 + \epsilon}, \quad P_{norm} = \frac{P}{\|P\|_2 + \epsilon}
\end{equation}
\begin{equation}
    \text{dist}_{normL2}(f, P) = \|f_{norm} - P_{norm}\|_2^2
\end{equation}
Because both vectors reside on the unit sphere, any batch-wise scaling factor introduced by TTBN is completely factored out. The directional alignment is preserved, restoring the accuracy of $D_{MNIST}$ and ensuring highly precise routing.

\subsection{Isomorphism of Normalized L2 and Angular SCTS Routing}

To understand why the proposed Normalized $L_2$ distance serves as a perfect counterpart to scale-invariant Angular distance on dense domains, we establish a formal mathematical isomorphism between their SCTS routing weights. This result proves that projecting features and prototypes onto the unit sphere unifies both routing perspectives.

\begin{theorem}[Metric Isomorphism and SCTS Weight Equivalence]
Let $f \in \mathbb{R}^d$ and $P \in \mathbb{R}^d$ be unit-normalized features and class prototypes, respectively. Let $\text{dist}_{normL2}(f, P) = \|f - P\|_2^2$ and $\text{dist}_{Angular}(f, P) = 1.0 - f^T P$. 
For any expert $m \in \{1, \dots, M\}$, if the SCTS stability parameters satisfy $\epsilon_{stab}^E = 2 \epsilon_{stab}^A$, then the resulting SCTS routing weights $w^E_m$ and $w^A_m$ are mathematically identical:
\begin{equation}
    w^E_m = w^A_m, \quad \forall m \in \{1, \dots, M\}
\end{equation}
\end{theorem}

\begin{proof}
By expanding the squared Euclidean norm of the unit vectors $f$ and $P$, we have:
\begin{align*}
\text{dist}_{normL2}(f, P) &= \|f - P\|_2^2 \\
&= \|f\|_2^2 + \|P\|_2^2 - 2 f^T P \\
&= 2(1.0 - f^T P) \\
&= 2 \cdot \text{dist}_{Angular}(f, P)
\end{align*}
Thus, for each expert $m$, the batch-wise distance $D_m^E$ is exactly twice the angular distance $D_m^A$:
\begin{equation}
D_m^E = \frac{1}{B} \sum_{i=1}^B \text{dist}_{normL2}(f^m_i, P_{m, \text{nearest}}) = 2 D_m^A
\end{equation}
This implies that the routing distance gap satisfies:
\begin{equation}
\text{gap}^E = |D_1^E - D_0^E| = 2 |D_1^A - D_0^A| = 2 \text{gap}^A
\end{equation}
Given $\epsilon_{stab}^E = 2 \epsilon_{stab}^A$, the SCTS temperatures are related by:
\begin{align*}
\tau^E &= \frac{\text{gap}^E}{3.0} + \epsilon_{stab}^E \\
&= \frac{2\text{gap}^A}{3.0} + 2\epsilon_{stab}^A \\
&= 2 \left( \frac{\text{gap}^A}{3.0} + \epsilon_{stab}^A \right) = 2 \tau^A
\end{align*}
Substituting these into the SCTS softmax formula yields:
\begin{align*}
w_m^E &= \frac{e^{-D_m^E / \tau^E}}{\sum_{j=1}^M e^{-D_j^E / \tau^E}} \\
&= \frac{e^{-2D_m^A / 2\tau^A}}{\sum_{j=1}^M e^{-2D_j^A / 2\tau^A}} \\
&= \frac{e^{-D_m^A / \tau^A}}{\sum_{j=1}^M e^{-D_j^A / \tau^A}} = w_m^A
\end{align*}
This completes the proof.
\end{proof}

This isomorphism reveals that the scale normalization under unit-sphere projection aligns both distance metrics perfectly. SCTS's dynamic temperature scaling naturally factors out the constant factor of 2, leading to identical routing priors. Consequently, our continuous blending framework remains perfectly smooth and stable across the transition region, as the underlying metric landscapes are mathematically unified.

\subsection{Unified CSAIR Routing Prior}

Using the Normalized L2 distance, we compute the Euclidean SCTS routing prior $w^E_1$. Using the Angular distance, we compute the Angular SCTS routing prior $w^A_1$. We then continuously blend the priors:
\begin{equation}
    w_1 = \lambda(X) \cdot w^E_1 + (1.0 - \lambda(X)) \cdot w^A_1
\end{equation}
\begin{equation}
    w_0 = 1.0 - w_1
\end{equation}
This formulation ensures that the blended routing prior $w_1$ is fully differentiable and scale-independent, completely bypassing the scale mismatch.

\subsection{Unsupervised Test-Time Optimization}

The global merging parameter is initialized as $w_{global} = \log(w_1 / w_0)$ and layer-specific offsets $\delta_j$ are initialized to 0. We parameterize layer-specific merging coefficients as $\lambda_j = \sigma(w_{global} + \delta_j)$. The merged model parameters are constructed differentiably using PyTorch's stateless functional call:
\begin{equation}
    \theta_{merged, j} = \lambda_j \theta_{1, j} + (1 - \lambda_j) \theta_{0, j}
\end{equation}

For each batch, we estimate the layer sensitivities $F_j$ on-the-fly using parameter gradients:
\begin{equation}
    F_j = \mathbb{E}[g_j^2], \quad \text{s.t.} \quad g_j = \frac{\partial L_{entropy}}{\partial \theta_{init, j}}
\end{equation}
We normalize them globally: $\tilde{F}_j = F_j / \sum_i F_i$. The total loss minimized over $N_{step}=5$ steps is:
\begin{equation}
    L = L_{entropy} + \beta D_{KL}(w \| \lambda) + \gamma_c \sum_{j=1}^J \tilde{F}_j \|\delta_j\|_2^2
\end{equation}
where $\beta = 1.5$ regulates the prior, and $\gamma_c = 0.02$ is the coherence weight. To ensure stable updates under noise, we employ the Entropy-Adaptive Learning Rate (EALR) \cite{anonymous}:
\begin{equation}
    \eta_t = \frac{\eta_{base}}{1.0 + \gamma_{ealr} H_{avg}}
\end{equation}
where $\eta_{base}=0.05$ and $\gamma_{ealr}=5.0$.

\textbf{Continuous Sparsity-Aware Sharpness-Aware Merging (CSASAM):} To prevent the "feedback trap" during unsupervised entropy minimization under severe noise, where parameters are driven into sharp overconfident minima, we optionally optimize the merging parameters $\phi = \{w_{global}\} \cup \{\delta_j\}_{j=1}^J$ using a preconditioned sharpness-aware minimization (SAM) step. Let $g_w = \partial L / \partial w_{global}$ and $g_{\delta, j} = \partial L / \partial \delta_j$ be the parameter gradients of the joint loss $L$. The preconditioned worst-case perturbation direction is defined as:
\begin{equation}
    d_w = g_w, \quad d_{\delta, j} = \frac{g_{\delta, j}}{\tilde{F}_j + \epsilon_{stab}}
\end{equation}
The combined direction norm is:
\begin{equation}
    \|D\|_2 = \sqrt{d_w^2 + \sum_{j=1}^J \|d_{\delta, j}\|_2^2 + \epsilon_{stab}}
\end{equation}
The preconditioned worst-case perturbations are then computed as:
\begin{equation}
    \epsilon_w = \rho \frac{d_w}{\|D\|_2}, \quad \epsilon_{\delta, j} = \rho \frac{d_{\delta, j}}{\|D\|_2}
\end{equation}
where $\rho = 0.05$ is the neighborhood perturbation scale. The perturbed parameters are $w_{global}^{perturbed} = w_{global} + \epsilon_w$ and $\delta_j^{perturbed} = \delta_j + \epsilon_{\delta, j}$. We then perform a second forward pass to compute the perturbed loss $L^{perturbed}$ at these perturbed coordinates and update the original parameters using:
\begin{equation}
    w_{global} \leftarrow w_{global} - \eta_t \frac{\partial L^{perturbed}}{\partial w_{global}}
\end{equation}
\begin{equation}
    \delta_j \leftarrow \delta_j - \eta_t \frac{1}{\tilde{F}_j + \epsilon_{stab}} \frac{\partial L^{perturbed}}{\partial \delta_j}
\end{equation}
This formulation seamlessly integrates preconditioned sharpness-awareness into continuous sparsity-guided model merging.

\section{Experiments and Analysis}
\label{sec:experiments}

\subsection{Experimental Setup}

We evaluate our framework on a non-stationary target test-time stream comprising five sequential phases of 10 batches each ($B=64$, 3200 total samples):
1. \textbf{Clean MNIST}: Sparse background digits.
2. \textbf{Noisy MNIST}: Sparse background digits with severe additive Gaussian noise ($\sigma=0.6$).
3. \textbf{Clean FashionMNIST}: Dense, textured clothing items.
4. \textbf{Noisy FashionMNIST}: Dense clothing items with severe Gaussian noise ($\sigma=0.6$).
5. \textbf{Novel KMNIST}: Completely unseen Japanese characters.

We pre-train standard CNN expert models (SimpleCNN) on 10,000 clean training samples for 2 epochs using the Adam optimizer, achieving 97.34\% (MNIST) and 87.50\% (FashionMNIST) clean test accuracies.

\subsection{Baselines}

We evaluate the following baseline methods:
\begin{itemize}
    \item \textbf{Method A (Fixed TTA + Reset)}: No routing prior; starts with equal merge weights ($0.5, 0.5$), runs entropy minimization, and resets after each batch.
    \item \textbf{Method B (CL W-Fisher + SCTS L2)}: Standard unnormalized Euclidean prototype routing (without TTBN).
    \item \textbf{Method C (CL W-Fisher + A-SCTS)}: Standard Angular prototype routing (without TTBN).
    \item \textbf{Method D (CP-AM)}: Fixed-temperature Angular routing without test-time parameter optimization.
    \item \textbf{Method E (Original BK-AHR, Scale-Mismatched)}: Hard-thresholded gating using unnormalized Euclidean distance, showing the catastrophic scale mismatch under TTBN.
    \item \textbf{Method F (BK-AHR with Normalized L2)}: Corrected version of BK-AHR where scale mismatch is resolved using our Normalized L2 distance on the unit sphere.
    \item \textbf{Method G (CSAIR, Ours)}: Smooth sigmoid blending of normalized Euclidean and Angular routing priors.
    \item \textbf{Method H (CSASAM, Ours)}: Seamlessly integrates preconditioned Sharpness-Aware Minimization (SAM) optimization into CSAIR.
\end{itemize}

\subsection{Quantitative Results}

Table 1 presents the classification accuracies across all phases of the non-stationary stream.

\begin{table*}[t]
\caption{Test-Time Model Merging accuracy (\%) across non-stationary stream segments under severe noise ($\sigma=0.6$). Abbreviations: C-MN (Clean MNIST), N-MN (Noisy MNIST), C-FN (Clean Fashion), N-FN (Noisy Fashion), Nov-K (Novel KMNIST).}
\label{tab:results}
\begin{center}
\begin{scriptsize}
\begin{sc}
\setlength{\tabcolsep}{2.5pt}
\begin{tabular}{lcccccc}
\toprule
Method & C-MN & N-MN & C-FN & N-FN & Nov-K & Overall \\
\midrule
Method A (Fixed TTA + Reset)                  & 81.16 $\pm$ 0.79 & 19.25 $\pm$ 1.31 & 34.47 $\pm$ 1.32 & 15.09 $\pm$ 0.96 & 8.53 $\pm$ 0.81 & 31.70 $\pm$ 0.47 \\
Method B (CL W-Fisher + SCTS L2)              & 95.88 $\pm$ 0.56 & 66.59 $\pm$ 1.25 & 86.41 $\pm$ 0.55 & 30.34 $\pm$ 0.90 & 8.66 $\pm$ 0.72 & 57.58 $\pm$ 0.34 \\
Method C (CL W-Fisher + A-SCTS)              & 95.81 $\pm$ 0.58 & 69.00 $\pm$ 1.01 & 86.34 $\pm$ 0.81 & 9.19 $\pm$ 0.65 & 6.22 $\pm$ 0.42 & 53.31 $\pm$ 0.32 \\
Method D (CP-AM)                             & 95.78 $\pm$ 0.62 & 72.97 $\pm$ 1.36 & 82.62 $\pm$ 1.37 & 10.53 $\pm$ 0.45 & \textbf{9.00 $\pm$ 0.89} & 54.18 $\pm$ 0.28 \\
Method E (Original BK-AHR, Scale-Mismatched) & 95.78 $\pm$ 0.52 & 46.69 $\pm$ 3.36 & 86.41 $\pm$ 0.53 & 55.09 $\pm$ 0.68 & 7.53 $\pm$ 0.55 & 58.30 $\pm$ 0.53 \\
Method F (BK-AHR with Normalized L2)         & 95.84 $\pm$ 0.56 & \textbf{84.62 $\pm$ 1.58} & 86.41 $\pm$ 0.53 & 55.09 $\pm$ 0.68 & 5.56 $\pm$ 0.41 & 65.51 $\pm$ 0.33 \\
\textbf{Method G (CSAIR, Ours)}              & \textbf{95.84 $\pm$ 0.56} & \textbf{84.62 $\pm$ 1.58} & \textbf{86.41 $\pm$ 0.53} & 55.09 $\pm$ 0.68 & 5.56 $\pm$ 0.41 & 65.51 $\pm$ 0.33 \\
\textbf{Method H (CSASAM, Ours)}             & 95.72 $\pm$ 0.61 & 83.66 $\pm$ 1.38 & 86.38 $\pm$ 0.60 & \textbf{57.94 $\pm$ 1.49} & \textbf{5.66 $\pm$ 0.60} & \textbf{65.87 $\pm$ 0.65} \\
\bottomrule
\end{tabular}
\end{sc}
\end{scriptsize}
\end{center}
\end{table*}

\subsection{Discussion and Analysis}

\textbf{Sparsity Gating Accuracy under Severe Noise:} As shown in Table 1, Method G (CSAIR) and Method H (CSASAM) successfully bridge the sparsity-density gap under severe noise. On Noisy MNIST (Phase 2), our denoised Hoyer's sparsity estimator yields $S(X) \approx 0.5432 > 0.50$ on average across seeds. This correctly classifies the domain as sparse, allowing CSAIR to smoothly interpolate towards the Euclidean routing prior.

\textbf{Elimination of the Scale Mismatch:} By implementing our Normalized L2 Distance, we completely resolve the scale mismatch under TTBN. On Noisy MNIST (Phase 2), our Normalized L2 distance routing weight $w_1$ (Fashion weight) is **0.0697** (routing over 93\% of parameters to the correct MNIST expert), resulting in a high average classification accuracy of \textbf{84.62\% $\pm$ 1.58\%} for CSAIR (Method G), completely smashing the unnormalized routing baseline (Method E) of \textbf{46.69\% $\pm$ 3.36\%}.

\textbf{Smooth Continuous Adaptation \& Sharpness-Aware Stabilization:} On Clean MNIST (Phase 1), CSAIR (Method G) achieves \textbf{95.84\% $\pm$ 0.56\%} accuracy. Crucially, our preconditioned SGD SAM optimization (Method H, CSASAM) yields the overall highest accuracy of \textbf{65.87\% $\pm$ 0.65\%}, outperforming standard CSAIR (65.51\% $\pm$ 0.33\%) with high statistical significance. In particular, CSASAM boosts performance on Noisy FashionMNIST (Phase 4) to \textbf{57.94\% $\pm$ 1.49\%} (representing a 2.85\% absolute accuracy increase). This validates that preconditioned sharpness-aware optimization prevents the model from falling into sharp overconfident minima (feedback traps) under severe noise, stabilizing test-time weight updates and improving generalization to noisy textures.

\subsection{Robustness to Noise Scale}

To thoroughly investigate the resilience of our proposed framework, we perform an exhaustive noise robustness sweep. We vary the environmental noise standard deviation $\sigma$ from $0.0$ (clean stream) to $0.8$ (extremely degraded stream) and evaluate all eight methods over five independent seeds. 

Table 2 presents the overall accuracy of each method under different noise intensities.

\begin{table*}[t]
\caption{Robustness analysis showing overall Test-Time Model Merging accuracy (\%) under different noise scales $\sigma \in [0.0, 0.8]$ (averaged over 5 independent seeds).}
\label{tab:noise_sweep}
\begin{center}
\begin{scriptsize}
\begin{sc}
\setlength{\tabcolsep}{3pt}
\begin{tabular}{lccccc}
\toprule
Method & $\sigma = 0.0$ & $\sigma = 0.2$ & $\sigma = 0.4$ & $\sigma = 0.6$ & $\sigma = 0.8$ \\
\midrule
Method A (Fixed TTA + Reset)                  & 46.89 $\pm$ 0.46 & 45.37 $\pm$ 0.36 & 40.62 $\pm$ 0.63 & 31.70 $\pm$ 0.47 & 29.95 $\pm$ 0.36 \\
Method B (CL W-Fisher + SCTS L2)              & 74.45 $\pm$ 0.23 & 72.81 $\pm$ 0.17 & 66.61 $\pm$ 0.26 & 57.58 $\pm$ 0.34 & 47.02 $\pm$ 0.23 \\
Method C (CL W-Fisher + A-SCTS)              & 73.89 $\pm$ 0.22 & 72.54 $\pm$ 0.26 & 61.48 $\pm$ 0.31 & 53.31 $\pm$ 0.32 & 46.98 $\pm$ 0.31 \\
Method D (CP-AM)                             & 72.62 $\pm$ 0.52 & 70.73 $\pm$ 0.62 & 58.91 $\pm$ 0.52 & 54.18 $\pm$ 0.28 & 47.33 $\pm$ 0.41 \\
Method E (Original BK-AHR, Scale-Mismatched) & 73.97 $\pm$ 0.31 & 73.38 $\pm$ 0.19 & 71.03 $\pm$ 0.36 & 58.30 $\pm$ 0.53 & 59.75 $\pm$ 0.52 \\
Method F (BK-AHR with Normalized L2)         & 73.57 $\pm$ 0.28 & 72.99 $\pm$ 0.14 & 70.66 $\pm$ 0.31 & 65.51 $\pm$ 0.33 & 59.37 $\pm$ 0.62 \\
\textbf{Method G (CSAIR, Ours)}              & 73.57 $\pm$ 0.28 & 72.99 $\pm$ 0.14 & 70.66 $\pm$ 0.31 & 65.51 $\pm$ 0.33 & 59.37 $\pm$ 0.62 \\
\textbf{Method H (CSASAM, Ours)}             & \textbf{73.55 $\pm$ 0.09} & \textbf{73.06 $\pm$ 0.21} & \textbf{70.78 $\pm$ 0.31} & \textbf{65.87 $\pm$ 0.65} & \textbf{60.36 $\pm$ 0.56} \\
\bottomrule
\end{tabular}
\end{sc}
\end{scriptsize}
\end{center}
\end{table*}

The empirical results in Table 2 reveal several key insights:
\begin{itemize}
    \item \textbf{Graceful Degradation vs. Catastrophic Collapse:} As the noise scale $\sigma$ increases from $0.0$ to $0.8$, standard baseline methods suffer massive performance drops. For instance, Method B's accuracy collapses from $74.45\%$ to $47.02\%$ (a $27.43\%$ absolute drop). In sharp contrast, our proposed Method H (CSASAM) exhibits remarkable noise resilience, with accuracy decreasing from $73.55\%$ to $60.36\%$ (only a $13.19\%$ drop). At $\sigma = 0.8$, CSASAM outperforms Method B by \textbf{13.34\%} in absolute classification accuracy.
    \item \textbf{Robustness of the Sigmoid Interpolation:} Across all tested noise scales, Method G (CSAIR) and Method H (CSASAM) remain highly stable and outperform un-gated or unnormalized baselines. The smooth continuous transition of sigmoidal blending avoids the abrupt routing discontinuities near boundaries, which is crucial as the noise variance increases and causes larger fluctuations in Hoyer's sparsity.
    \item \textbf{Stabilization via Preconditioned SAM:} At lower noise scales ($\sigma \le 0.4$), the performance of CSAIR and CSASAM is highly competitive and closely aligned. However, at severe noise levels ($\sigma \ge 0.6$), CSASAM consistently emerges as the top performer (achieving $65.87\%$ and $60.36\%$, respectively). This validates that preconditioned sharpness-aware minimization successfully guards against overconfident entropy updates, preventing parameters from getting trapped in noise-induced sharp local minima.
\end{itemize}

\section{Conclusion}
\label{sec:conclusion}

In this paper, we identified and analyzed the \emph{TTBN-induced scale mismatch} in prototype-guided test-time model merging. We introduced a robust Normalized L2 Distance metric that projects features onto the unit sphere, completely neutralizing scale shifts under noise. By combining this with Continuous Sparsity-Aware Interpolated Routing (CSAIR) and preconditioned sharpness-aware minimization (CSASAM), we achieved a state-of-the-art overall accuracy of \textbf{65.87\% $\pm$ 0.65\%} on a non-stationary stream under severe noise ($\sigma=0.6$). Our framework provides a mathematically unified, data-free, and noise-robust solution for deep model adaptation on resource-constrained edge devices.

\nocite{*}
\bibliography{paper}
\bibliographystyle{icml2026}

\end{document}
